% Transcription — fonds Alexandre Grothendieck, shelfmark 62, batch 3 (pages 41-60).
% Established from the facsimile archives/batches/62/batch-03.pdf (a copy of
% 62.pdf fetched through the project relay, no cover sheet: PDF page k =
% archivists' page 40+k), per the skill transcribe-grothendieck. The folder is
% ninety-six pages in five batches; this is the third. Batch 1's notation is
% kept (underlined sheaves \underline{L}, \underline{O}_S; the Lie sheaf as
% \mathcal{L}; E_n as \mathcal{E}_n; « ? » for an unreadable index in a
% formula). His gothic S is set \mathfrak{S}_3, as elsewhere in the fonds.
%
% Pass: Opus 5.5 (claude-opus-5-5), 2026-09-26 — first pass, unchecked against the pages by a human.
%
% Pages skipped: 44, 48, 50, 52, 54 — leaves of the same unrelated
% typescript as batch 1's page 19 and batch 2's pp. 21-40 (« n° 276 »,
% typed pp. 16, 13 upside down, 27, 26, 25: complétés, anneaux de
% Zariski), the other sides of his manuscript sheets, with nothing of his
% on them; 55, a blank sheet. Page 42 (typed p. 17) carries one pencilled
% marginal word, « topologiques », called into the text by a stroke; as
% batch 2 does for its pencilled typed pp. 21 and 23, it is kept with one
% \note, the attribution to him flagged as uncertain. Page 46 is a leaf of that typescript (its p. 15)
% carrying a few ink jottings of his, written sideways in the margin
% ((0,1,∞), (X', X''), ...), which belong to the S_3 discussion of
% pp. 41-49: it gets one \note and his jottings, nothing of the typed text.
% Page 58 is a tan cover sheet inscribed in pencil top right « Théorie
% transcendante (Généralités classiques) »: no \page marker, its words become
% a section heading with a \note.
%
% Pages summarised: none. Pages 56 and 57 (blue paper) are scratch sheets
% of isolated formulas, given in page order.
%
% His pagination: none visible on these leaves. Batch 2's numbered-formula
% run 1)-(12) (pp. 27-39) does not continue here: no numbered formula on
% pp. 41-53. Numbered runs: on p. 60,
% « Théorème II », « Théorème III », a theorem whose numeral is an ink blot,
% « Théorème IV », « Théorème V » (Théorème I is not labelled on the page:
% a cross stands before « Si une fonction elliptique n'a pas de pôles... »).
% No date on any page.
%
% What the batch is about. Pages 41-53 continue batch 1's description of an
% elliptic curve in residue characteristics ≠ 2,3 as a double cover of a
% Brauer-Severi scheme: X a Brauer-Severi scheme (conic) over S with a
% non-ramified trisection X' (p. 41); such data are classified by
% H^1(S, S_3), X' gives an « associated » X''; in char. ≠ 2,3 a unique
% isomorphism X -> P^1_S sending X', X'', X_1 to 0, ∞, 1 (p. 43); with a
% section s disjoint from X' one gets the invariant j, a section of O_S; if
% j and j-1 are units, the inverse image of j : S -> P^1_S is the Galois
% S_3-cover defining X; j = 0 reduces the structure group to Z/2 (H^1(S,Z/2)),
% j = 1 to Z/3 (H^1(S,Z/3)); the factorisation X -> P^1_S -> P^1_S,
% t -> j t^2 (pp. 45-47); the description by an S_3-torsor S' and
% t in Γ(S', O_S') with t, t-1 units and σ.t = φ_σ(t) (λt = 1/t,
% μt = 1-t), « fonctoriel », and the example S = Spec(O), O complete local
% with algebraically closed residue field (p. 49); « Autre façon de
% procéder » (p. 51): X - s_∞ = Λ a principal homogeneous bundle under the
% line bundle L, the trisection giving a barycentric section, hence
% Λ ≅ L and the data become an invertible sheaf L with a non-ramified
% trisection of V(L) of sum zero (p. 53). Page 56: scratch on the
% Weierstrass form (D_ξ u_ξ)^2 = 4u_ξ^3 + b_ξ u_ξ + c_ξ, its behaviour under
% ξ -> λξ (weights 2, 4, 6), the flag E_1 ⊂ E_2 ⊂ ... and powers of a sheaf
% written as a triangle. Page 57: a flatness/descent computation with
% B ⊗_A B and Nakayama. Pages 59-60 (brown ink, a different, earlier-looking
% fair hand, under the cover of p. 58): classical elliptic functions —
% definition, period parallelograms, Théorèmes I-V (Liouville, residues,
% order, zeros/poles, sum of zeros, algebraic relation).
%
% Hardest pages: 41, 45, 51, 53 (a very fast abbreviated hand: formulas
% carry, connective prose often does not); the slanted, partly struck margin
% note on 53. Counts: 238 \ill{}, 41 \uncertain{}. Readings to check first:
% « sur X » at the end of the trisection's definition (p. 41, one expects
% S); « t ↦ j t^2 » (p. 47); the exponents in « Δ^n ≃ L^18 » (p. 56).

\documentclass[11pt,a4paper]{article}
% Le préambule commun à toutes les transcriptions du fonds Grothendieck.
%
% Il tient en une idée : **l'incertitude se compose, elle ne se résout pas.**
% Une transcription qui lisse un mot douteux a détruit la seule chose pour
% laquelle on la faisait — un lecteur doit pouvoir savoir, à chaque endroit,
% ce qui était sur le feuillet et ce qui a été ajouté. Les six macros qui
% suivent sont donc l'appareil critique, et non de la décoration.
%
% Les mêmes commandes sont reconnues par `scripts/render.mjs`, qui produit la
% vue de lecture du site. Une macro ajoutée ici doit l'être aussi là-bas,
% faute de quoi le rendu échouera bruyamment — ce qui est voulu : un rendu
% silencieusement faux est pire qu'un rendu absent.

\ProvidesPackage{grothendieck}[2026/08/08 Transcriptions du fonds Grothendieck]

\usepackage{amsmath,amssymb,amsthm}
\usepackage{mathrsfs}
\usepackage{stmaryrd}
% Les diagrammes commutatifs sont partout dans ce fonds : c'est la langue
% même de la géométrie algébrique de Grothendieck, et une transcription qui
% les rendrait en prose ne transcrirait rien.
\usepackage{tikz-cd}
% Français par défaut — c'est la langue des feuillets — mais l'anglais est
% chargé aussi : la traduction et le résumé sont des documents anglais, et la
% typographie française (espaces avant les deux-points) y serait une faute.
% Ces documents basculent par \selectlanguage{english} en tête de corps.
\usepackage[english,french]{babel}
\usepackage{fontspec}
\usepackage{geometry}
\geometry{a4paper,margin=25mm,marginparwidth=28mm}
\usepackage{marginnote}
\usepackage{xcolor}
% ulem, not soul. `soul`'s \st and \ul are built on a character-by-character
% scan that XeLaTeX's font machinery breaks: the document fails with an
% undefined control sequence rather than degrading. ulem does the same two jobs
% and is Unicode-engine safe. `normalem` stops it hijacking \emph, which the
% transcriptions use for Grothendieck's own emphasis.
\usepackage[normalem]{ulem}
\usepackage{hyperref}
\hypersetup{colorlinks=true,linkcolor=black,urlcolor=black,pdfborder={0 0 0}}

% --- Métadonnées du lot -----------------------------------------------------
% Reprises telles quelles par le script de rendu, qui les affiche en tête de
% la vue de lecture. Elles fixent ce que le fichier prétend couvrir : sans
% elles, une transcription flotte sans point d'attache dans un fonds de seize
% mille feuillets.

\newcommand{\folder}[1]{\gdef\grothFolder{#1}}
\newcommand{\batch}[1]{\gdef\grothBatch{#1}}
\newcommand{\leaves}[2]{\gdef\grothFirst{#1}\gdef\grothLast{#2}}
\newcommand{\foldertitle}[1]{\gdef\grothFoldertitle{#1}}
\gdef\grothFolder{?}\gdef\grothBatch{?}\gdef\grothFirst{?}\gdef\grothLast{?}\gdef\grothFoldertitle{}

% --- Le résumé --------------------------------------------------------------
%
% Il ouvre la lecture modernisée, sous le titre « Résumé », et il est composé
% en petit. Ce n'est pas un ornement : c'est le seul passage du document qui
% ne soit pas des mathématiques, et un lecteur doit pouvoir le reconnaître
% avant de l'avoir lu — puis le sauter s'il connaît déjà le sujet, ou s'y
% arrêter s'il ne le connaît pas. Le filet qui le suit marque la reprise du
% corps.

\newenvironment{resume}
  {\small\setlength{\parskip}{0.45em}}
  {\par\medskip\hrule\medskip}

% --- Mots-clés ---------------------------------------------------------------
%
% En anglais, en clôture du résumé de la lecture modernisée : le vocabulaire
% moderne sous lequel chercher ce que les feuillets construisent. Le manifeste
% du site (`npm run manifest`) les extrait pour étiqueter la cote entière —
% cette ligne est la source unique des tags, il n'y a pas de fichier à côté.
\newcommand{\keywords}[1]{\par\smallskip\noindent
  {\footnotesize\textsc{Keywords} --- \textit{#1}}\par}

% --- L'appareil critique ----------------------------------------------------

% Le numéro de feuillet, en marge. C'est l'ancre qui permet de régler un
% désaccord sur une formule en regardant *un* feuillet plutôt qu'une chemise
% de six cents. Le site s'en sert aussi pour tourner les pages du fac-similé
% quand on fait défiler la transcription.
\newcommand{\leaf}[1]{%
  \marginnote{\footnotesize\color{gray!70}\textsf{#1}}%
  \label{leaf:#1}\ignorespaces}

% Illisible : on ne devine pas. Un mot inventé qui se lit comme les autres est
% le pire résultat possible.
\newcommand{\ill}{\textcolor{red!60!black}{[\,\ldots\,]}}

% Lecture douteuse : le mot est donné, et signalé comme incertain.
\newcommand{\uncertain}[1]{\uline{#1}}

% Ajout de l'éditeur — une lettre manquante, un mot sous-entendu.
\newcommand{\add}[1]{\textcolor{blue!50!black}{[#1]}}

% Biffé par Grothendieck lui-même. Ce qu'il a rayé fait partie du document :
% une rature dit souvent où la pensée a changé de direction.
\newcommand{\struck}[1]{\sout{#1}}

% Note du transcripteur, sur l'état matériel du feuillet.
\newcommand{\note}[1]{\par\smallskip{\small\color{gray!60!black}\textit{[#1]}}\par\smallskip}

% Note marginale de Grothendieck, distincte du corps du texte.
\newcommand{\marginal}[1]{\par\smallskip\noindent\hspace{1em}%
  {\small\color{gray!60!black}#1}\par\smallskip}

% --- Titre ------------------------------------------------------------------

\AtBeginDocument{%
  \begin{center}
    {\large\bfseries Fonds Alexandre Grothendieck --- cote n\textsuperscript{o}~\grothFolder}\\[2pt]
    {\itshape\grothFoldertitle}\\[4pt]
    {\small feuillets \grothFirst--\grothLast\ (lot \grothBatch)}\\[2pt]
    {\footnotesize\color{gray!60!black}Transcription établie d'après le fac-similé numérisé
     par l'Université de Montpellier.\\
     Les lectures incertaines sont soulignées, les illisibles notées [\,\ldots\,],
     les ajouts entre crochets.}
  \end{center}
  \vspace{1em}}

\endinput


\folder{62}
\batch{3}
\pages{41}{60}
\dating{[à partir de 1963]}
\watermark{Édition de démonstration}
\foldertitle{Courbes elliptiques (généralités - vieille rédaction) : notes manuscrites (s.d.)}

\begin{document}

\page{41}
Soit $S$ un schéma de base, et soit $X$ un schéma de Br.-Severi sur $S$ muni d'une
\emph{trisection} non ramifiée $X'$, i.e.\ un sous-schéma \add{\uncertain{fini}}
de $X$ qui est un revêtement \uncertain{étale} sur $X$, de degré $3$ sur $X$.
\note{ainsi sur la page : « sur $X$ » les deux fois ; on attendrait $S$. Le mot
en interligne, en haut à droite, se lit mal}
Un tel truc peut être considéré comme un \ill{} objet : \ill{} structure
$\mathfrak{S}_3$, puisque \ill{} $\mathbf{P}^1_S$ \ill{} \ill{} il a été dit. La
classification des tels objets \ill{} \ill{} \ill{} $=$ \ill{} la classification
des \ill{} de $H^1(S,\mathfrak{S}_3)$ des \struck{\ill{}} \add{\uncertain{gal.}}
\emph{\uncertain{rev.\ non ramifiés}} \ill{} groupe $\mathfrak{S}_3$ de $S$. De plus,
on définit \ill{} $X$ une \ill{} \ill{}\add{\uncertain{involutive}}, \ill{} \ill{}
\ill{} \ill{} $X'$ en un $X''$, dit \emph{associé} à $X'$.
\note{le mot souligné est lu « associé » ; le mot écrit au-dessus de la ligne et
appelé par un trait est douteux}
Évidemment, \ill{} \ill{} \ill{} \ill{} $(X,X'')$, \ill{} \ill{} \ill{} \ill{}, et
\ill{} \ill{} $X'$ [\ill{} \ill{} \ldots]. De plus, on \uncertain{obtient}
[\ill{} \uncertain{descente} \ldots] une \emph{bisection}, \ill{}
\uncertain{ramifiée} \struck{[\ill{} \ill{}} \ill{} \ill{} \ill{} \ill{} $S$ \ill{}
car.\ $\neq 2,3$] \ldots. \ill{} \ill{} \ill{} \ill{} \ill{}, \struck{\ill{}} \ill{}
$\mathfrak{S}_3$, \ill{} \ill{} \ill{} \ill{} :
\[
\mathbf{P}^1_S \xrightarrow{\ p\ } \mathbf{P}^1_S
\]
(\ill{} \ill{} $\mathfrak{S}_3$ \ill{}

\page{42}
\note{page 17 du même tapuscrit « n° 276 » (complété d'un groupe topologique
abélien, Lemme 1 sur les suites exactes de groupes métrisables), sans rapport
direct avec les courbes elliptiques. Il porte au crayon, dans la marge gauche,
en face de « suite de groupes abéliens », un mot appelé par un trait dans le
texte, d'une main qui paraît être la sienne, sans certitude}
\marginal{\uncertain{topologiques}}

\page{43}
\ill{} \ill{} \ill{} \ill{} \ill{} \ill{} \ill{}) \ill{} \uncertain{isomorphisme}
\uncertain{unique} [\ill{} $x$ \emph{\uncertain{pt} de car.\ $\neq 2,3$}]
\[
X \xrightarrow{\ f\ } \mathbf{P}^1_S
\]
qui est \uncertain{caractérisé} \ill{} les sections $0$, $\infty$, $1$
\add{\ill{} de $\infty$ dans $X'$} \struck{\ill{}} exclusivement : l'image
\ill{} \ill{} \ill{} \ill{} \ill{}, celle de $\infty$ est $X''$, celle de $1$ est
$X_1$.
\note{l'ajout est écrit sous « $0,\infty,1$ », entre deux traits ; le début de la
phrase suivante, « l'image \ldots », se lit mal. Après cette phrase, une croix
isolée sépare les deux parties de la page}

\uncertain{Supposons} maintenant qu'on se donne, en plus, une donnée
\uncertain{supplémentaire}, une section $s$ de $X$ \struck{\ill{}} sur $S$,
\ill{} \ill{} par $X'$ [i.e.\ une section de $X-X'$ sur $S$]. Son image par $f$ est
une section de $\mathbf{P}^1_S-(\infty)$, i.e.\ une section de $\underline{O}_S$,
appelée l'\emph{invariant} de $(X,X',s)$, soit $j$.
\note{dans $(X,X',s)$ une lettre est biffée devant le $s$}
\uncertain{Supposons} par \ill{} que l'invariant $j$ ne prenne pas les valeurs
$0$, $1$, i.e.\ soit inversible, ainsi que $j-1$. Alors l'image inverse dans
$\mathbf{P}^1_S$ \struck{\ill{}} de $j\colon S\to\mathbf{P}^1_S$ définit un
\emph{\struck{\ill{}} \add{revêtement} galoisien de $S$, non ramifié, de groupe
$\mathfrak{S}_3$}, \struck{\ill{}} et c'est \emph{celui qui définit $X$}. [À
l'inverse des faisceaux

\page{45}
catégoriques généraux \ldots].

Dans ce cas, l'invariant suffit à reconstituer \struck{\ill{}} la structure
$(X,X',s)$.

Dans le cas où $X$ est \ill{}, \ill{}, et où \emph{$j$ n'est pas identiquement
égal à $0$ ou $1$}, \ill{} \ill{} \ill{} \ill{} \ill{} \struck{\ill{}} \ill{}
\ill{} \ill{} $X$ \ill{} \ill{} $U$ \ill{} c'est \ill{} \ill{} [$\mathbf{P}^1_S$
\ill{}], \ill{} \ill{} \ill{} \ill{}, d'où reconstitution de $(X,X')$ ; \ill{}
\ill{} \ill{}, on \ill{} $s$ \ill{} $U$, donc ($X$ \ill{}) partout.
\note{le passage souligné sous « $j$ n'est pas identiquement égal à $0$ ou $1$ »
est souligné en deux traits séparés}
Dans \ill{} \ill{}, il faut des $H$ \emph{\uncertain{supplémentaires}}. \ill{}
\uncertain{en particulier}, \ill{} \ill{} \ill{} \ill{} : \uncertain{quelle} \ill{}
\uncertain{convenablement} ??
Il faut exprimer (i) qu'une certaine extension de corps \ill{} de degré $3$
\ill{} du corps des fonctions de $X$ est non ramifiée sur $X$ (ii) qu'une
certaine application \uncertain{rationnelle} de $X$ qui se déduit \ill{}
\uncertain{définie} \ldots

\page{46}
\note{feuillet du tapuscrit sans rapport (n° 276, p. 15 : corollaires sur les
anneaux noethériens) ; seules sont données les notations qu'il y a jetées à
l'encre dans la marge droite, écrites de biais}
\marginal{$(2,\ \tfrac12(2-1))$ \quad $(x',\ x'')$ \quad $(0,1,\infty)$}
\note{chacune est entourée d'une accolade ; sur le tapuscrit, $X_1$ est barré
d'une croix à l'encre dans « $K[[X_1,\ldots,X_r]]$ »}

\page{47}
\struck{\uncertain{Supposons} \ill{} \ill{} \ill{} \ill{} \ill{}}

En tous cas, si $j=0$, la situation est décrite par le fait que l'on a fait une
réduction de \ill{} structural de $\mathfrak{S}_3$ à $\mathbf{Z}/2$
\add{[$X''$ \struck{\ill{} $X'$} \ill{} une section \ill{} !]}, et \ill{} \ill{}
est déterminée par la donnée d'un élément de $H^1(S,\mathbf{Z}/2)$.
\struck{\ill{} \ill{}} \add{Et, si $j=1$}, \ill{} la situation est décrite
\ill{} \ill{} \ill{} \ill{} une réduction du groupe structural :
$\mathbf{Z}/3\subset\mathfrak{S}_3$ [\ill{} : une section \ill{}
\add{la bisection} $X_1$] et la situation est décrite par la donnée d'un élément
de $H^1(S,\mathbf{Z}/3)$.
\note{« Et, si $j=1$ » est écrit au-dessus d'un mot biffé ; « la bisection » est
en interligne au-dessus de $X_1$}

\struck{\ill{}} En tous les cas, si $j$ est inversible, alors faisant opérer
$\mathfrak{S}_3$ dans $\mathbf{P}^1_S/(\mathbf{Z}/3)$ \ill{} \ill{} \ill{}
\uncertain{critique}, \ill{} \ill{} dans \struck{\ill{}} \add{\ill{}
\ill{}} \add{(\ill{} de $X'$, $X''$ et $s$)} trois \emph{sections} distinctes,
\ill{} \ill{} \ill{} \ill{} \ill{} $\infty$, $0$, $1$, d'où une identification de
ce dernier à $\mathbf{P}^1_S$ ; et $X\to\mathbf{P}^1_S$ se factorise par
\[
\begin{array}{ccccc}
X & \longrightarrow & \mathbf{P}^1_S & \longrightarrow & \mathbf{P}^1_S \\
 & & \infty,0,1 & \longmapsto & \infty,0,j
\end{array}
\qquad t\mapsto jt^2 ;
\]
la bisection naturelle \ill{} \ill{} s'identifie donc \ill{} : $\sqrt{j}$,
\note{le $S$ du premier $\mathbf{P}^1_S$ est surchargé ; « $t\mapsto jt^2$ » est
écrit au-dessous de la flèche, la lettre lue $j$ sous réserve}

\page{49}
\struck{\ill{}} \uncertain{façon}

Ceci nous a permis de donner de façon très explicite \ill{} \ill{} la
« courbe \add{\ill{}} » de \uncertain{genre} $0$ sur $S$, munie d'une section et
d'une trisection non ramifiée, disjointes. [\struck{\ill{}} $=$ revêtement
\add{$\mathfrak{S}_3$} \struck{\ill{}} galoisien \add{$S'$} de groupe
$\mathfrak{S}_3$ de $S$ \struck{\ill{}} \uncertain{pour} une section $t$ \struck{\ill{}}
disjointe de $0,1,\infty$, \emph{invariante}
\struck{\ill{}}
\[
S'\times_{(S,\mathfrak{S}_3)}\mathbf{P}^1_S
\]
par le groupe $\mathfrak{S}_3$, i.e.\ \struck{une section}
\[
t\in\Gamma(S',\underline{O}_{S'}),
\]
inversible et telle que $t-1$ soit inversible, et satisfaisant à
\[
\sigma.t = \varphi_\sigma(t) \qquad \sigma\in\mathfrak{S}_3
\]
i.e.
\[
\begin{cases} \lambda t = \dfrac{1}{t} \\[1ex] \mu t = 1-t \end{cases}
\]
\note{« invariante » est écrit sous une accolade qui embrasse « une section $t$
disjointe de $0,1,\infty$ » ; devant $S'\times_{(S,\mathfrak{S}_3)}\mathbf{P}^1_S$,
deux essais sont biffés. Le $\varphi_\sigma$ est écrit sur une autre lettre}

\begin{tikzcd}
X & \mathbf{P}^1_{S'} \arrow[l] \arrow[d] \\
S & S' \arrow[l, "\mathfrak{S}_3"]
\end{tikzcd}

\note{carré en marge gauche ; $X\to S$ est tracé d'une flèche fine, et un long
trait vertical, portant en biais « fonctoriel », borde le passage depuis le
crochet ouvrant}
\marginal{fonctoriel}

\emph{Exemple} $S=\operatorname{Spec}(\mathcal{O})$, $\mathcal{O}$ anneau
local \emph{complet} \add{(de car.\ $\neq 2,3$)} à corps résiduel alg.\ clos
\struck{\ill{}}. Alors \struck{\ill{} \ill{}} $\pi_1(S)=0$, donc la situation
équivaut, \add{\uncertain{Dém.} \ill{} \ill{}} \ill{} à la donnée d'un
$t\in\mathcal{O}^*$ tel que \struck{\ill{}} $t(t-1)\in\mathcal{O}^*$, deux
tels $t$, $t'$ donnant des \uncertain{objets} \uncertain{isomorphes}
\struck{\ill{}} s.s.s.\ $t'$ \ill{} de $t$ par \ill{} $\varphi_\sigma(t)$,
$\sigma\in\mathfrak{S}_3$.
\marginal{non fonctoriel}
\note{« non fonctoriel », en biais dans la marge gauche, est souligné deux
fois ; l'ajout « Dém.\ \ldots » est entouré au-dessus de la ligne}

\section*{Autre façon de procéder}

\note{titre de sa main, souligné, en tête de la page 51}

\page{51}
Prenons la section privilégiée comme section $\infty$, donc
$X-s_\infty=\Lambda$ apparaît comme un $S$-fibré principal \struck{projectif}
homogène de groupe $\mathcal{L}$, où $\mathcal{L}$ est un \ill{} \ill{} \ill{},
associé à un \struck{fibré} faisceau inversible $\underline{L}$ [\ldots\ \ill{}
explique de façon bien précise.
\note{un double trait vertical, dans la marge gauche, borde les lignes qui
suivent}
\ill{} schéma de Brauer--Severi et une section dans le dernier \ldots]. Alors la
\emph{trisection} de $\Lambda$ permet de construire, \ill{} \ill{} [\ldots] une
section de $\Lambda^3$, d'où [$3$ \uncertain{étant} inversible] une section de
$\Lambda$, i.e.\ la \emph{section barycentrique} [\ill{} explique \ill{} \ill{}
\ill{}, \ill{} \ill{} \ill{} \ill{} \ill{} inversible \ill{} un schéma \ill{}
\add{\uncertain{défini}}, \ill{} $\mathcal{L}$ \ill{} \ill{} \ill{} une
\ill{}-section \ill{} \ill{} fibré principal $\Lambda$ de groupe $\mathcal{L}$].
Cela définit un isomorphisme
\[
\Lambda \simeq \mathcal{L}
\]
\emph{global}, identifiant $X'$ à une

\page{53}
trisection de $\mathcal{L}$, de « somme » nulle.

\emph{La \struck{classification} \add{donnée} d'un système} \struck{privilégié}
\emph{\ill{} \add{portés} sur $S$} \ill{} \ill{} \uncertain{équivalente}
\emph{\ill{}} à \emph{celle d'un faisceau} \struck{\ill{}} \emph{inversible}
\struck{\ill{}} $\underline{L}$ \add{\emph{non ramifiée}} \emph{sur $S$, et d'une
trisection} $X'$ \emph{de} $\mathcal{L}=V(\underline{L})$, \emph{de}
\struck{\ill{}} \emph{« somme » nulle} \struck{\ill{}}
\struck{$V(\underline{L})^{\otimes 3}$}.
\note{la phrase est soulignée par morceaux ; « non ramifiée », entouré, est
écrit sous « inversible $\underline{L}$ » et relié par un trait à « $X'$ »}
Mais on peut aussi éliminer les \add{\ill{}} \ill{}, symétriques \ill{}
\struck{\ill{}} \ill{} trisection.
\marginal{\ill{} \ill{} \ill{} \ill{} \ill{} \ill{} \ill{} \ill{}}
\note{note en biais dans la marge gauche, sur une dizaine de lignes courtes,
en partie biffée d'un trait oblique ; on y distingue $\mathfrak{S}_3$ et une
croix, sans pouvoir la lire}

\page{56}
\note{feuille de brouillon (papier bleuté) : formules isolées, données dans
l'ordre de la page ; un petit rectangle griffonné en haut à gauche, portant
$\Delta^2$, $\Delta^3$, $y$, $x$ aux coins, n'est pas repris}

\[
(D_\xi u_\xi)^2 = 4u_\xi^3 + b_\xi u_\xi + c_\xi
\]
\struck{$(D_\xi x_\xi)^2 = 4\ldots$} \quad \struck{$u_\xi\in\Delta^?$}
\[
(D_{\lambda\xi}u_{\lambda\xi})^2 = 4u_{\lambda\xi}^3 + b_{\lambda\xi}u_{\lambda\xi}
+ c_{\lambda\xi}
\]
\[
u_\xi = 6x_{\xi^2} \qquad\qquad u_{\lambda\xi} = \rho\, u_\xi
\]
\note{une lettre biffée devant $u_{\lambda\xi}=\rho u_\xi$}
\[
C \,\big|\, \underline{O}\to\underline{O}(1) \qquad
f_*(\underline{O}(n)) = E_n \qquad\Big|\qquad
\lambda^2\rho^2(D_\xi u_\xi)^2 = 4\rho^3u_\xi^3 + \rho b_{\lambda\xi}u_\xi
+ c_{\lambda\xi}
\]
\note{$C$ au-dessus d'une flèche verticale vers $S$ ; sous le dernier membre,
un signe d'égalité vertical}
\[
\underline{O} = E_1 \subset E_2 \subset E_3 \subset E_4 \subset E_5 \subset E_6
\subset \cdots \subset E_n
\]
\[
\underline{O} = \Delta^0 \qquad L_2 \qquad L_3 \qquad\Big|\qquad L_2^2 \qquad
L_2L_3 \qquad L_2^3
\]
\note{sous chaque inclusion une petite accolade, qui en désigne le quotient ;
sous $L_3$, « $\simeq \Delta L_2$ » ; sous $L_2^2$, $L_2L_3$, $L_2^3$, des
termes en $\Delta$ ($\Delta L_3$, $\Delta^3 L_2$, \ldots) couverts de hachures ;
sous $L_2$ trois lignes illisibles, où l'on lit $\omega_1,\omega_2,\omega_3$, et
sous le tout une accolade marquée « \uncertain{pairs} »}
\[
y^2 = 4x^3 + bx + c \qquad\qquad \frac{y}{x} = D
\]
\note{une lettre biffée devant $y^2$, et sous « $\frac{y}{x}$ » le $x$ est
barré}
\[
(D_\xi u_\xi)^2 = 4\,\frac{\rho}{\lambda^2}\,u_\xi^3 +
\frac{1}{\rho\lambda^2}\,b_{\lambda\xi}u_\xi + \frac{c_{\lambda\xi}}{\lambda^2\rho^2}
\]
\[
\rho = \lambda^2 \qquad
\begin{cases}
u_{\lambda\xi} = \lambda^2 u_\xi\\
b_{\lambda\xi} = \lambda^4 b_\xi\\
c_{\lambda\xi} = \lambda^6 c_\xi
\end{cases}
\qquad
\begin{array}{l}
u_\xi \in L_2 \simeq \Delta^2\\
b_\xi \in \Delta^0\\
c_\xi \in \Delta^0
\end{array}
\]
\note{ainsi sur la page : $\Delta^0$ pour $b_\xi$ et $c_\xi$ ; l'exposant de
$b_\xi$ est surchargé}
\[
\Delta^4\longrightarrow\underline{O}_S \qquad \Delta^6\longrightarrow\underline{O}_S
\qquad \Delta^2\longrightarrow\Delta^2 \qquad\qquad
\begin{array}{l} \Delta^{?}\simeq\mathcal{L}^{18}\\ \Delta^2\simeq\mathcal{L}^3 \end{array}
\]
\note{les deux relations de droite, plus petites, se lisent mal ; l'exposant de
$\mathcal{L}$ est lu $18$ sous réserve}
\[
\begin{cases}
b\in\Delta^{-4}\\
c\in\Delta^{-6}\\
u \ \text{scalaire}
\end{cases}
\qquad\qquad u_\xi = \xi^2
\]
\note{sous « $u_\xi=\xi^2$ », un trait horizontal prolongé d'une oblique, et
$\xi^2$ au-dessous}

\page{57}
\note{feuille de brouillon (papier bleuté), dans l'ordre de la page}

\[
\begin{array}{c} X\\[1ex] \\ S \end{array} \qquad
X\times_S X \supset R \qquad\qquad
\begin{array}{ccc}
\mathcal{B} & & A\\
X & \longrightarrow & Y\\
\nearrow & & \nearrow\\
X\times_Y X & \longrightarrow & X
\end{array}
\qquad x_i :
\]
\[
B\otimes_A B
\]
\[
\begin{array}{c} B\\ | \\ A \end{array} \qquad
A \longrightarrow B\otimes_A B \qquad B
\]
$x_i\otimes 1$ base de $B\otimes B$ sur $B\otimes B$ \quad / \ill{} \ill{} $\mathfrak{m}\subset A$
\note{les $\otimes$ de cette ligne sont surchargés ; « $\mathfrak{m}\subset A$ »
est lu sous réserve}

$x_i\otimes 1$ \quad --- \quad $L\otimes_K L$ sur $1\otimes L$

$x_i$ base de $L$ sur $k$

$x_i$ générateurs de $B$ sur $A$ par Nakayama.
\[
\sum\lambda_i x_i = 0 \qquad \lambda_i\in A
\]
\struck{$\sum\lambda_i x_i$}
\[
\Bigl(\sum\lambda_i x_i\Bigr)\otimes_A 1 =
\sum(x_i\otimes 1)(1\otimes\lambda_i) = 0
\]
$1\otimes\lambda_i=0$ pour \uncertain{tout} $i$ ;

\uncertain{Utilisant} $B\otimes_A B\to B$, on trouve
\struck{$\lambda_i\otimes 1 = 0$ \quad $1\otimes\lambda_i = \lambda_i\otimes 1$}
\[
\lambda_i = 0 \ \text{dans}\ B
\]

\section*{Théorie transcendante (Généralités classiques)}

\note{inscrit au crayon, de sa main, en haut à droite du feuillet de garde
(page 58), qui ne porte rien d'autre ; le feuillet ne reçoit pas de numéro de
page. Les pages 59--60, sur papier brun, sont d'une encre brune et d'une
écriture posée, plus ancienne d'aspect}

\page{59}
\emph{Fonction elliptique.} C'est une fonction méromorphe dans le plan à
distance finie, et admettant deux périodes distinctes $\omega_1$ et $\omega_2$.
\struck{\ill{}} En particulier, une constante est une fonction elliptique. Si
$f(z)$ n'est pas une constante, le rapport de $\omega_1$ et $\omega_2$ n'est pas
réel \struck{imaginaire}.
\add{Définition de la période ($\omega_1$ et $\omega_2$ donnent \ill{} \ill{}
\ill{} \uncertain{parallélogramme}). Mais il faut remplacer « $\omega_1$ et
$\omega_2$ » !}
\add{le domaine d'une fonction elliptique n'est \uncertain{certain} \ill{} \ill{}
\ill{}}
\note{deux lignes serrées en interligne, au-dessus de « \ldots\ correspondant du
plan » ; lues par fragments}
\ill{} correspondant du plan. On peut définir un réseau de parallélogrammes
correspondants aux périodes $\omega_1$ et $\omega_2$ et à un sommet $a$ arbitraire.
Ils sont déterminés par leurs sommets de la forme $a+n_2\omega_2+n_1\omega_1$ ;
tout point qui n'est pas de la forme $a+n\omega_2+\lambda\omega_1$, ou
$a+\lambda\omega_2+n\omega_1$ ($\lambda$ et $n$ réels, $n$ entier) appartient : $1$
et $1$ seul de ces parallélogrammes \add{il \uncertain{bien} \ill{}} \struck{\ill{}},
tout point qui est de la forme parallèle, mais non de la forme
$a+n_1\omega_2+n_1\omega_1$ ($=a+\omega$), \struck{\ill{}} appartient : deux et
deux seulement de ces par., il appartient aux côtés, non au sommet ; enfin, un
point de la forme $a+n_1\omega_2+n_2\omega_1$ appartient : $4$ et $4$ seulement de
tels par., il appartient : leurs sommets.
\note{ainsi sur la page : $n_1$ deux fois dans la deuxième forme ; les indices
varient d'une forme à l'autre}

Si on se donne un parallélogramme de périodes, tout point est l'homologue soit
d'un pt intérieur $\xi$, soit de deux pts sur les côtés, et non les sommets,
soit de $4$ sommets.
\add{En particulier, deux pts intérieurs d'un même par.\ de périodes ne sont pas
homologues}
\struck{Si on se donne la fonction}

On peut trouver des par.\ des périodes, tels que sur leurs côtés ne se trouve
aucun zéro, pôle etc. ; dans ce cas, la \struck{\ill{}} \add{liste} des zéros
(pôles) avec leur ordre de multiplicité, est indépendante du par.\ considéré. Il
en est ainsi en toute généralité, si on \uncertain{compte} une seule fois
\struck{\ill{}} les zéros de tout groupe de zéros homologues
(\struck{situés} ils ont situés sur les sommets --- auquel cas on les divise par
$4$ --- autrement où les divise par deux).

\page{60}
L'ens.\ des val.\ prises par $f(z)$ dans un par.\ des périodes est égal à
l'ensemble des valeurs prises dans tout le plan.

$\times$ Si une fonction elliptique n'a pas de pôles, elle est constante.
\note{la croix est dans la marge, devant la phrase}

Remarque que la somme, le produit, le quotient, la dérivée, de fonctions
elliptiques de m.\ périodes sont de m.\ périodes.

\emph{Théorème II} La somme des résidus d'une fonction elliptique dans un par.\
de périodes est nulle.

Donc une fonction elliptique \ill{} admet dans un par.\ de périodes au moins
deux pôles (chacun compté avec son ordre de multiplicité).

\emph{Théorème III} Le nombre des zéros et le \uncertain{nombre} des pôles d'une
fonction elliptique dans un parallélogramme des périodes est le même (en
supposant toujours que sur les côtés \ldots) ; il en est de même du nombre de
fois qu'il prend \struck{la valeur} \add{une valeur} donnée (sous les m.\
réserves).

Le nombre est appelé ordre de la fonction elliptique.

\emph{Théorème \ill{}} Deux fonctions elliptiques de m.\ périodes, qui ont mêmes
zéros et mêmes pôles, avec le m.\ ordre de multiplicité, sont égales à un
facteur constant près. Elles sont égales à une constante additive près, si elles
ont m.\ périodes, et m.\ pôles, avec m.\ partie principale.
\note{le chiffre romain de ce théorème est recouvert d'une tache d'encre ; un
long trait courbe, dans la marge gauche, le relie au théorème II}

\emph{Théorème IV} La somme des zéros d'une fonction \struck{toute} elliptique
dans un parallélogramme de périodes égale la somme de ses pôles, à une période
près.

\emph{Théorème V.} Entre deux fonctions elliptiques de m.\ périodes existe une
relation algébrique.

\end{document}
